\section*{Chapter \ref{ch:rs:trade-flow-features} --- Trade-Flow Features}

\begin{solution}{exo:rs:trade-flow-features:1}
The mid is 20.005. \omterm{def:rs:trade-flow-features:sign}{Tick rule}: unclassified (no earlier price), $+1$, $+1$ (a zero tick keeps the last sign), $-1$, $-1$. Lee--Ready: every
trade is off the mid, so it is the \omterm{def:rs:trade-flow-features:sign}{quote rule}: $-1$, $+1$, $+1$, $-1$, $-1$. The two differ on the first trade only.
\end{solution}

\begin{solution}{exo:rs:trade-flow-features:2}
$\Phi(0.8) = 0.788$: 78.8\% of the \omterm{def:rs:market-data-for-research:bar}{bar}'s volume is classified as buys, whatever the trades inside it actually were.
\end{solution}

\begin{solution}{exo:rs:trade-flow-features:3}
Decay exponent $\gamma = \alpha - 1 = 0.4$; Hurst exponent $H = (3 - \alpha)/2 = 0.8$.
\end{solution}

\begin{solution}{exo:rs:trade-flow-features:4}
The imbalances $|V^B - V^S| = |2V^B - 1\,000|$ are 400, 200, 800, 0 and 600, summing to 2\,000; \omterm{def:rs:trade-flow-features:vpin}{VPIN} $= 2\,000/(5 \times 1\,000) = 0.4$.
\end{solution}

\begin{solution}{exo:rs:trade-flow-features:5}
$30 \times 0.10 = 3.0$ ticks. Signed trade volume sees only the executions; the quotes also move when orders are cancelled or added, which
OFI counts and signed volume does not. On the simulated tape the $R^2$ is 14\% against OFI's 53\%, so the slope is estimated from a
relation that leaves most of the price change unexplained.
\end{solution}

\begin{solution}{exo:rs:trade-flow-features:6}
Take a Poisson process with intensity $\lambda_1$ for half the time and $\lambda_2 > \lambda_1$ for the other half, each regime long. No event
causes another, but an event is more likely to fall in the busy regime, and then more events follow it soon: counts are over-dispersed
compared with one Poisson process. A Hawkes process with a constant baseline has only one way to produce that clustering, excitation, so the
fit puts mass in the kernel, and the more the regimes differ the higher the branching ratio. On the chapter's tape the fit gives 0.65
against the 0.4 planted, which the flat tape recovers (0.403).
\end{solution}

\begin{solution}{exo:rs:trade-flow-features:7}
\texttt{rs\_tradeflow.stale\_quotes}: 99.7\% with quotes 0.05 seconds late, 97.9\% at 0.5 seconds and 93.8\% at 2 seconds. The loss is about
five points per second of staleness at first and slows (6.2 points at two seconds): the older the quote, the more likely the mid has moved
past the trade price, but a quote that has moved once may move back.
\end{solution}

\begin{solution}{exo:rs:trade-flow-features:8}
It is a statement chosen after the event. A warning is useful only with its false-alarm rate: how often \omterm{def:rs:trade-flow-features:vpin}{VPIN} reached such levels with no crash,
and whether high readings were followed by high volatility on average (Andersen and Bondarenko found they were a poor \omterm{def:rs:anatomy-of-a-predictor:predictor}{predictor} of it). A
highest reading of the year is also one draw among the year's hours; some hour before any event is the maximum of its neighbourhood. And the
level depends on choices (bucket size, window, the signing method) that must be fixed before the event for the claim to count.
\end{solution}

\begin{solution}{pb:rs:trade-flow-features:1}
\begin{enumerate}
\item 15\% before the episode, 45\% during it: the informed share nearly triples.
\item $-0.36$ ticks before, $+0.96$ during (the mid ten seconds after the trade, minus the trade price, in the direction of the trade).
\item Aggressors pay the half-spread, half a tick, and most of them are noise traders whose orders carry no information: the resting side
earns the spread.
\item Yes: the aggressors win and the liquidity providers lose on average, the definition of toxic flow.
\item Buckets of 2\,649 shares, $1/400$ of the session's volume, so 400 buckets, and a window of 20 buckets.
\item 0.31 before, 0.35 during: a rise of 13\%.
\item At minute 119.7, 20 minutes after the episode ends.
\item 16 seconds after the episode starts; \omterm{def:rs:trade-flow-features:vpin}{VPIN} is above the same threshold for 6.0\% of the session outside the episode.
\item The informed traders buy when the efficient price is above the mid and sell when it is below; as that price wanders both ways, their flow
over a bucket (27 seconds on average) is not one-sided. The one-sided runs come from the noise traders' \omterm{def:rs:trade-flow-features:acf}{metaorders}, which the episode replaces
with informed flow rather than adds to.
\item $-0.38$ before, $0.84$ during; it peaks at 5\,740 seconds, inside the episode.
\item At 5\,440 seconds, 40 seconds after the start: each value waits ten seconds for its mark-out, and the rolling one-minute mean needs
toxic trades to fill it.
\item 8.4\%.
\item For studying toxicity after the fact, the mark-out: it is toxicity by definition. To protect a market maker in real time, the mark-outs
of its own recent fills at a short horizon together with the order-book features of chapter~8; \omterm{def:rs:trade-flow-features:vpin}{VPIN} has not shown here that it would help.
\item Informed traders with a signal that points one way for longer than a bucket, a single large informed order split over minutes, or news
that moves the value in one step. \omterm{def:rs:trade-flow-features:vpin}{VPIN} measures that case, not this one.
\item 0.65 on this tape; 0.403 on the flat tape, against the planted 0.4.
\item 0.10 ticks per 100 shares of net buying, with an $R^2$ of 14\%.
\item \omterm{def:rs:trade-flow-features:vpin}{VPIN} crosses its threshold first, 16 seconds in, at 6.0\% false alarms; the mark-out 40 seconds in, at 8.4\%. Neither is clean, and
\omterm{def:rs:trade-flow-features:vpin}{VPIN}'s early crossing is not matched by a rise in its level.
\item \textbf{Named result.} When the informed share of the flow goes from 15\% to 45\% and the aggressors' mark-out turns from $-0.36$ to
$+0.96$ ticks, \omterm{def:rs:trade-flow-features:vpin}{VPIN} rises only 13\%, from 0.31 to 0.35, and reaches its maximum 20 minutes after the episode ends.
\item A toxicity measure must be tested against toxicity itself (the mark-outs), on data where it is known, before it is trusted where it
is not; a plausible story linking the measure to informed trading is not a test.
\item Toxic flow is order flow whose aggressors are informed, so that the liquidity providers who trade against it lose on average.
\end{enumerate}
\end{solution}

\begin{solution}{iq:rs:trade-flow-features:1}
With quotes, compare the trade price with the prevailing mid, taking the quotes as of the trade's timestamp and checking how the two clocks
align (Lee and Ready, 1991); fall back on the \omterm{def:rs:trade-flow-features:sign}{tick rule} at the mid. Without quotes, the \omterm{def:rs:trade-flow-features:sign}{tick rule}. Measure the error where the truth is
known (a feed with an aggressor flag, or a simulator): on the chapter's tape a quote one second stale loses 3.7 points of accuracy.
\end{solution}

\begin{solution}{iq:rs:trade-flow-features:2}
\omterm{def:rs:trade-flow-features:acf}{Metaorders} are split into many child orders of one sign, so signs persist (Lillo, Mike and Farmer, 2005). Prices stay close to unpredictable
because liquidity adjusts: as a run of buys continues, the impact of each further buy falls, and sellers refill the ask, so the predictable sign
does not become a predictable price move of the same size.
\end{solution}

\begin{solution}{iq:rs:trade-flow-features:3}
Flow is toxic when those who trade against it lose on average. Measure it by mark-outs: for each of the book's fills, the mid a few horizons
later minus the fill price, from the book's side, by client, venue, time of day and size; a mark-out that turns negative quickly is toxic
flow.
\end{solution}

\begin{solution}{iq:rs:trade-flow-features:4}
A baseline that varies (time of day, news, volatility regimes) produces clusters that a constant-baseline fit attributes to excitation. Fit
with a baseline that follows the time of day or a measured activity level, compare on held-out data, and check on simulated data with a
known branching ratio: the chapter's tape gives 0.65, its flat version 0.403 against 0.4.
\end{solution}

\begin{solution}{iq:rs:trade-flow-features:5}
Cut trades into equal-volume buckets, sign them (or classify \omterm{def:rs:market-data-for-research:bar}{bars} by BVC), and average the absolute buy--sell imbalance over the last $n$
buckets, divided by the bucket size; the claim is that one-sided flow is toxic. The criticism: it was a poor \omterm{def:rs:anatomy-of-a-predictor:predictor}{predictor} of short-run volatility
(Andersen and Bondarenko, 2014), its level depends on arbitrary choices, and one-sided flow need not be informed, nor informed flow
one-sided.
\end{solution}

\begin{solution}{iq:rs:trade-flow-features:6}
If \omterm{def:rs:trade-flow-features:acf}{metaorder} sizes have tail exponent $\alpha$, the sign autocorrelation decays as $\tau^{-(\alpha-1)}$ (Lillo, Mike and Farmer, 2005), and a
series whose autocorrelation decays as $\tau^{-\gamma}$ with $0 < \gamma < 1$ has partial sums with variance growing like $n^{2 - \gamma}$,
so $H = 1 - \gamma/2 = (3 - \alpha)/2$. Test: aggregate the executions into orders, compute the variance of block sums of signs across scales
and regress in logs. The chapter's tape, with $\alpha = 1.5$, predicts 0.75 and gives 0.715.
\end{solution}
